\documentclass[12pt]{article}

\usepackage[margin=1in]{geometry}
\usepackage{enumitem}
\usepackage{amsmath}
\usepackage{array}
\usepackage{titlesec}

\title{Lecture 2 Notes: \\ Aristotle's \textit{Categories}}
\author{}
\date{}

\begin{document}

\maketitle

\section*{Central Theme}

Aristotle's \textit{Categories} teaches us how to sort what we say about things
by asking what kind of being, predicate, or expression we are dealing with.

\[
\boxed{
\textit{Categories} = \text{a grammar of being and predication}
}
\]

Before Aristotle studies arguments, demonstrations, syllogisms, or scientific
knowledge, he begins with the basic ways things can be named, said, and
classified.

\section*{1. Why Begin with Categories?}

The \textit{Categories} is one of Aristotle's first logical works. It does not
begin with arguments. It begins with terms.

This is important. Aristotle thinks that before we can judge whether an
argument is valid, we need to understand what kinds of things our words signify.

\subsection*{Teaching Point}

Logic begins with speech. But speech must first be clarified. Aristotle begins
by asking:

\begin{quote}
When we say something about something, what kind of thing are we saying?
\end{quote}

\subsection*{Whiteboard}

\[
\text{Words} \rightarrow \text{Terms} \rightarrow \text{Predicates}
\rightarrow \text{Assertions} \rightarrow \text{Arguments}
\]

\section*{2. Equivocal, Univocal, and Derivative Naming}

Aristotle begins by distinguishing three ways names are used.

\subsection*{Equivocal Naming}

Things are named equivocally when they have the same name, but the definition
corresponding to that name is different.

For example, a real man and a picture of a man may both be called ``animal,''
but they are not animals in the same sense.

\[
\text{Equivocal} = \text{same name, different definition}
\]

\subsection*{Univocal Naming}

Things are named univocally when they share both the same name and the same
definition.

For example, a man and an ox are both called ``animal'' in the same sense,
because the definition of animal applies to both.

\[
\text{Univocal} = \text{same name, same definition}
\]

\subsection*{Derivative Naming}

Things are named derivatively when their name comes from another name, but with
a different ending.

For example:

\[
\text{grammar} \rightarrow \text{grammarian}
\]

\[
\text{courage} \rightarrow \text{courageous}
\]

\subsection*{Teaching Point}

This opening distinction matters because Aristotle wants to discipline our
language. A shared word does not always mean a shared essence.

\subsection*{Whiteboard}

\[
\boxed{
\text{Same name} \neq \text{same definition}
}
\]

\section*{3. Simple and Composite Expressions}

Aristotle next distinguishes simple expressions from composite expressions.

\subsection*{Simple Expressions}

Simple expressions are single terms, such as:

\[
\text{man}, \quad \text{ox}, \quad \text{runs}, \quad \text{wins}
\]

These do not, by themselves, assert anything true or false.

\subsection*{Composite Expressions}

Composite expressions combine terms, such as:

\[
\text{the man runs}
\]

\[
\text{the man wins}
\]

These can be true or false because they make an assertion.

\subsection*{Teaching Point}

Truth and falsity do not belong to isolated terms. They belong to assertions.

\subsection*{Whiteboard}

\[
\text{``man''} \quad \text{not true or false}
\]

\[
\text{``the man runs''} \quad \text{true or false}
\]

\[
\boxed{
\text{Term} \neq \text{Assertion}
}
\]

\section*{4. Predicable of a Subject and Present in a Subject}

This is one of the most important distinctions in the \textit{Categories}.

Aristotle distinguishes between things that are said \textbf{of} a subject and
things that are present \textbf{in} a subject.

\subsection*{Predicable of a Subject}

Something is predicable of a subject when it can be said of that subject.

For example:

\[
\text{man is predicable of an individual man}
\]

\[
\text{animal is predicable of man}
\]

\subsection*{Present in a Subject}

Something is present in a subject when it exists in that subject and cannot
exist separately from it.

For example:

\[
\text{a particular whiteness is present in a body}
\]

\[
\text{a particular grammatical knowledge is present in a mind}
\]

\subsection*{The Fourfold Distinction}

\[
\begin{array}{c|c|c}
 & \textbf{Predicable of a Subject} & \textbf{Not Predicable of a Subject} \\
\hline
\textbf{Present in a Subject} &
\text{knowledge} &
\text{this whiteness} \\
\hline
\textbf{Not Present in a Subject} &
\text{man} &
\text{this individual man} \\
\end{array}
\]

\subsection*{Teaching Point}

This table is the conceptual center of the lecture. Aristotle is sorting the
ways things relate to subjects.

Some things are said \textbf{of} subjects. Some things exist \textbf{in}
subjects. Some do both. Some do neither.

\subsection*{Whiteboard}

\[
\text{said of} \neq \text{present in}
\]

\[
\boxed{
\text{The individual substance is neither said of nor present in a subject.}
}
\]

\section*{5. The Ten Categories}

Aristotle then lists the basic kinds of simple expression.

\[
\begin{array}{ll}
1. & \text{Substance} \\
2. & \text{Quantity} \\
3. & \text{Quality} \\
4. & \text{Relation} \\
5. & \text{Place} \\
6. & \text{Time} \\
7. & \text{Position} \\
8. & \text{State} \\
9. & \text{Action} \\
10. & \text{Affection}
\end{array}
\]

\subsection*{Examples}

\[
\begin{array}{ll}
\textbf{Substance} & \text{man, horse, Socrates} \\
\textbf{Quantity} & \text{two cubits long, three cubits long} \\
\textbf{Quality} & \text{white, grammatical, courageous} \\
\textbf{Relation} & \text{double, half, greater} \\
\textbf{Place} & \text{in the marketplace, in the Lyceum} \\
\textbf{Time} & \text{yesterday, last year} \\
\textbf{Position} & \text{lying, sitting} \\
\textbf{State} & \text{shod, armed} \\
\textbf{Action} & \text{to lance, to cauterize} \\
\textbf{Affection} & \text{to be lanced, to be cauterized}
\end{array}
\]

\subsection*{Teaching Point}

The categories are not merely a vocabulary list. They are Aristotle's attempt
to classify the basic kinds of things that can be said.

The guiding question is:

\begin{quote}
What kind of predicate is this?
\end{quote}

\subsection*{Classroom Example}

Take the sentence:

\begin{quote}
Socrates is sitting in Athens today, wearing sandals, teaching Plato's students.
\end{quote}

We can sort its elements categorically:

\[
\begin{array}{ll}
\text{Socrates} & \text{substance} \\
\text{sitting} & \text{position} \\
\text{in Athens} & \text{place} \\
\text{today} & \text{time} \\
\text{wearing sandals} & \text{state} \\
\text{teaching} & \text{action}
\end{array}
\]

\section*{6. Substance: The Central Category}

Substance is the most important category in the work.

Aristotle distinguishes between primary substance and secondary substance.

\subsection*{Primary Substance}

Primary substance is the individual thing.

Examples:

\[
\text{this individual man}
\]

\[
\text{this individual horse}
\]

Primary substance is neither predicable of a subject nor present in a subject.

\[
\boxed{
\text{Primary substance is the basic subject.}
}
\]

\subsection*{Secondary Substance}

Secondary substances are the species and genera to which primary substances
belong.

Examples:

\[
\text{this man} \rightarrow \text{man} \rightarrow \text{animal}
\]

Here:

\[
\text{this man} = \text{primary substance}
\]

\[
\text{man} = \text{species}
\]

\[
\text{animal} = \text{genus}
\]

Species and genera are called secondary substances because they tell us what
the primary substance is.

\subsection*{Species and Genus}

Aristotle says that species are more truly substance than genera because they
give a more precise account of the individual.

To say:

\[
\text{Socrates is a man}
\]

is more informative than to say:

\[
\text{Socrates is an animal}
\]

Both are true, but the species gives a more determinate account.

\subsection*{Teaching Point}

Primary substance anchors the whole system. Everything else is either said of
primary substances or present in primary substances.

\subsection*{Whiteboard}

\[
\boxed{
\text{If primary substances did not exist, nothing else could exist.}
}
\]

\[
\text{this man} \rightarrow \text{man} \rightarrow \text{animal}
\]

\section*{7. Marks of Substance}

Aristotle identifies several marks of substance.

\subsection*{Substance Has No Contrary}

There is no contrary of an individual man or an individual horse.

\[
\text{No contrary: this man}
\]

\subsection*{Substance Does Not Admit Degrees}

A man cannot be more or less man.

Someone can be more or less pale, more or less healthy, or more or less
learned, but not more or less human.

\[
\text{Quality admits degree; substance does not.}
\]

\subsection*{Substance Can Admit Contrary Qualities}

The same individual substance can receive contrary qualities at different
times.

For example:

\[
\text{Socrates is healthy}
\]

\[
\text{Socrates is sick}
\]

The same individual remains one and the same substance while undergoing change.

\subsection*{Teaching Point}

This is one of the most important metaphysical ideas in the lecture:

\[
\boxed{
\text{Substance remains the same while qualities change.}
}
\]

\section*{8. Quantity}

Quantity is divided into discrete and continuous quantity.

\subsection*{Discrete Quantity}

Discrete quantities have parts that do not share a common boundary.

Examples:

\[
\text{number}
\]

\[
\text{speech}
\]

\subsection*{Continuous Quantity}

Continuous quantities have parts that share a common boundary.

Examples:

\[
\text{line}, \quad \text{surface}, \quad \text{solid}, \quad \text{time}, \quad \text{place}
\]

\subsection*{The Mark of Quantity}

The distinctive mark of quantity is that equality and inequality can be
predicated of it.

\[
\boxed{
\text{Quantity is what can be equal or unequal.}
}
\]

\section*{9. Relation}

Relatives are things whose meaning refers to something else.

Examples:

\[
\text{double} = \text{double of something}
\]

\[
\text{half} = \text{half of something}
\]

\[
\text{master} = \text{master of a slave}
\]

\[
\text{knowledge} = \text{knowledge of something}
\]

\subsection*{Teaching Point}

A relative cannot be understood in isolation. Its meaning points beyond itself.

\subsection*{Whiteboard}

\[
\boxed{
\text{Relation} = \text{being understood by reference to another}
}
\]

\section*{10. Quality}

Quality answers the question:

\begin{quote}
What sort of thing is it?
\end{quote}

Aristotle gives several kinds of quality.

\subsection*{Habits and Dispositions}

A habit is more stable and lasting. A disposition is more easily changed.

Examples of habits:

\[
\text{knowledge}, \quad \text{virtue}
\]

Examples of dispositions:

\[
\text{being warm}, \quad \text{being ill}, \quad \text{being healthy}
\]

\subsection*{Capacities and Incapacities}

Some qualities are capacities or incapacities.

Examples:

\[
\text{good boxer}, \quad \text{good runner}, \quad \text{healthy}, \quad \text{sickly}
\]

\subsection*{Affective Qualities and Affections}

Some qualities affect perception.

Examples:

\[
\text{sweetness}, \quad \text{bitterness}, \quad \text{heat}, \quad \text{cold}
\]

\subsection*{Shape and Figure}

Shape also belongs to quality.

Examples:

\[
\text{triangular}, \quad \text{quadrangular}, \quad \text{straight}, \quad \text{curved}
\]

\subsection*{The Mark of Quality}

The distinctive mark of quality is likeness and unlikeness.

\[
\boxed{
\text{Quality is that in virtue of which things are like or unlike.}
}
\]

\section*{11. Action, Affection, Place, Time, Position, and State}

Aristotle treats some categories more briefly.

\subsection*{Action}

Action concerns what something does.

Examples:

\[
\text{heating}, \quad \text{cutting}, \quad \text{teaching}
\]

\subsection*{Affection}

Affection concerns what is done to something.

Examples:

\[
\text{being heated}, \quad \text{being cut}, \quad \text{being taught}
\]

\subsection*{Place}

Place answers:

\[
\text{Where?}
\]

Example:

\[
\text{in the Lyceum}
\]

\subsection*{Time}

Time answers:

\[
\text{When?}
\]

Example:

\[
\text{yesterday}
\]

\subsection*{Position}

Position concerns posture or arrangement.

Examples:

\[
\text{sitting}, \quad \text{lying}, \quad \text{standing}
\]

\subsection*{State}

State concerns a condition of being equipped or clothed.

Examples:

\[
\text{shod}, \quad \text{armed}
\]

\section*{12. Opposition}

Aristotle distinguishes four kinds of opposition.

\[
\begin{array}{ll}
1. & \text{Correlatives} \\
2. & \text{Contraries} \\
3. & \text{Positives and Privatives} \\
4. & \text{Affirmation and Negation}
\end{array}
\]

\subsection*{Correlatives}

Examples:

\[
\text{double} \leftrightarrow \text{half}
\]

\[
\text{master} \leftrightarrow \text{slave}
\]

\subsection*{Contraries}

Examples:

\[
\text{good} \leftrightarrow \text{bad}
\]

\[
\text{white} \leftrightarrow \text{black}
\]

Some contraries have intermediates. For example, between white and black there
may be grey.

\subsection*{Positives and Privatives}

Examples:

\[
\text{sight} \leftrightarrow \text{blindness}
\]

\[
\text{possession} \leftrightarrow \text{privation}
\]

Privation is not merely absence. It is the absence of something that naturally
should be present.

\subsection*{Affirmation and Negation}

Examples:

\[
\text{he sits}
\]

\[
\text{he does not sit}
\]

This kind of opposition is especially important because here one statement is
true and the other false.

\subsection*{Teaching Point}

Opposition is not one simple relation. Aristotle distinguishes different ways
things can be opposed.

\[
\boxed{
\text{Not every opposite is a contrary.}
}
\]

\section*{13. Priority, Simultaneity, Motion, and Having}

The final sections of the \textit{Categories} refine important philosophical
terms.

\subsection*{Priority}

Something can be prior in several senses:

\begin{itemize}[itemsep=0.25em]
    \item prior in time
    \item prior by dependence
    \item prior in arrangement
    \item prior in honor or importance
    \item prior as cause
\end{itemize}

\subsection*{Simultaneity}

Things can be simultaneous in time, or simultaneous by nature.

For example, double and half are simultaneous by nature because each implies the
other.

\subsection*{Motion}

Aristotle lists six kinds of motion:

\[
\begin{array}{ll}
1. & \text{generation} \\
2. & \text{destruction} \\
3. & \text{increase} \\
4. & \text{diminution} \\
5. & \text{alteration} \\
6. & \text{change of place}
\end{array}
\]

\subsection*{Having}

The term ``to have'' is used in many ways.

Examples:

\[
\text{to have knowledge}
\]

\[
\text{to have a height}
\]

\[
\text{to have a coat}
\]

\[
\text{to have a house}
\]

\section*{14. Aristotle's Method}

Throughout the \textit{Categories}, Aristotle follows a repeated method.

\[
\boxed{
\text{define the kind} \rightarrow \text{divide it} \rightarrow \text{identify its mark}
}
\]

For example:

\[
\begin{array}{lll}
\textbf{Category} & \textbf{Division} & \textbf{Distinctive Mark} \\
\hline
\text{Substance} & \text{primary / secondary} &
\text{admits contraries while remaining one} \\
\text{Quantity} & \text{discrete / continuous} &
\text{equal and unequal} \\
\text{Quality} & \text{habits, capacities, affections, shape} &
\text{like and unlike} \\
\text{Relation} & \text{various correlatives} &
\text{reference to another}
\end{array}
\]

\section*{15. Conclusion}

The \textit{Categories} is not merely a list of words. It is Aristotle's first
attempt to map the basic structures of meaningful speech about reality.

Before Aristotle can ask whether an argument is valid, he asks what our words
are doing when we say something about something.

\[
\boxed{
\text{To classify speech is already to begin classifying being.}
}
\]

\section*{Final Framing}

The lecture should leave students with three main ideas:

\begin{enumerate}[itemsep=0.5em]
    \item Aristotle begins logic by clarifying language.
    \item The most important distinction is between what is said of a subject
    and what is present in a subject.
    \item Primary substance is the anchor of Aristotle's categorical system.
\end{enumerate}

\[
\boxed{
\textit{Categories} = \text{Aristotle's first map of what can be said of what is.}
}
\]

\end{document}